\part{Aktueller Stand}

\chapter{Was zu tun ist}

\section{Frau B., ein letztes Mal}

\begin{merksatz}
\boxtitel{Was der Band aus einer Frage gemacht hat}
Am Anfang stand: Gerinnungshemmung oder nicht?

\begin{center}
{\small
\renewcommand{\arraystretch}{1.5}
\begin{tabular}{@{}P{1.2cm}P{3.4cm}P{3.6cm}@{}}
\toprule
\textbf{Kap.} & \textbf{Frage} & \textbf{Antwort}\\
\midrule
$1$ & Was rechnet man? & Erwartungsnutzen; Schwelle $1{,}5\,\%$.\\
$2$ & Woher der Nutzen? & Standardlotterie --- und er gehört ihr.\\
$3$ & Lohnt Abklärung? & Höchstens $1{,}5$ gesunde Tage.\\
$4$ & Lohnt ein Test? & Nur zwischen zwei Schwellen.\\
$5$ & Wie sicher ist das? & Kipppunkte, Zweiwegeprüfung.\\
$6$ & Und über Jahre? & $1{,}38$ QALY Gewinn.\\
$7$--$9$ & Und morgen wieder? & Rückwärts rechnen, erkunden.\\
$10$--$12$ & Was fehlt der Rechnung? & Kosten, Ambiguität, der Mensch.\\
$13$--$14$ & Und für alle? & Arrow, und die Grenzen der Summe.\\
\bottomrule
\end{tabular}}
\end{center}

\begin{center}
\begin{minipage}{0.9\textwidth}
\textbf{Die Empfehlung hat sich in fünfzehn Kapiteln nicht geändert. Was sich geändert hat, ist die Zahl der Stellen, an denen man sie bestreiten kann.}
\end{minipage}
\end{center}
\end{merksatz}

\section{Vier Muster}

\begin{merksatz}
\boxtitel{Muster 1 --- die Bewertung ist nie neutral}
\begin{center}
{\small
\renewcommand{\arraystretch}{1.48}
\begin{tabular}{@{}P{1.4cm}P{7.0cm}@{}}
\toprule
$1$ & Jede Leitlinienschwelle enthält ein Verhältnis von Schäden --- meist ungenannt.\\
$2$ & Standardlotterie, Zeitabwägung und Skala liefern für denselben Zustand verschiedene Zahlen.\\
$12$ & Dieselbe Zahl in vier Darstellungen erzeugt vier Entscheidungen.\\
$14$ & Die Summe der QALY enthält bereits eine Gerechtigkeitsvorstellung.\\
\bottomrule
\end{tabular}}
\end{center}

\begin{center}
\begin{minipage}{0.88\textwidth}
\kursiv{\textbf{Wer keine Bewertung angibt, hat trotzdem eine benutzt.} Das ist die erste Warnung des Vorworts, und sie gilt auf jeder Ebene --- beim einzelnen Patienten, in der Leitlinie, in der Erstattungsentscheidung. Der Gewinn der Rechnung ist nicht Objektivität, sondern Sichtbarkeit.}
\end{minipage}
\end{center}
\end{merksatz}

\begin{merksatz}
\boxtitel{Muster 2 --- der Wert liegt in der Änderung}
\begin{center}
{\small
\renewcommand{\arraystretch}{1.48}
\begin{tabular}{@{}P{1.4cm}P{7.0cm}@{}}
\toprule
$3$ & Ein Befund, der nichts ändert, ist wertlos --- gleich wie genau er ist.\\
$4$ & Unter- und oberhalb der Testschwellen ist jeder Test überflüssig.\\
$5$ & Eine Größe ist nur wichtig, wenn ihre Spanne die Entscheidung kippt.\\
$9$ & Erkunden lohnt nur, solange man die Kenntnis noch anwenden kann.\\
\bottomrule
\end{tabular}}
\end{center}

\begin{center}
\begin{minipage}{0.88\textwidth}
\kursiv{\textbf{Das ist der brauchbarste Gedanke des ganzen Bandes} --- und er lässt sich ohne jede Rechnung anwenden: Was täte ich bei welchem Ergebnis? Wären die Antworten gleich, ist die Untersuchung entbehrlich. Die Empfindlichkeitsanalyse aus Kapitel~5 ist dieselbe Frage, auf Modellgrößen angewandt.}
\end{minipage}
\end{center}
\end{merksatz}

\begin{merksatz}
\boxtitel{Muster 3 --- Zeit verändert alles}
\begin{center}
{\small
\renewcommand{\arraystretch}{1.48}
\begin{tabular}{@{}P{1.4cm}P{7.0cm}@{}}
\toprule
$6$ & Was in einem Jahr Tausendstel sind, wird in zwanzig Jahren mehr als ein Lebensjahr.\\
$6$ & Die Diskontierung entscheidet über Vorbeugung gegen Behandlung.\\
$7$ & Wer rückwärts rechnet, findet Reihenfolgen, die vorwärts nicht sichtbar sind.\\
$9$ & Beim letzten Patienten erkundet man nicht mehr.\\
\bottomrule
\end{tabular}}
\end{center}

\begin{center}
\begin{minipage}{0.88\textwidth}
\kursiv{\textbf{Die zweite Zeile ist politisch folgenreich.} Drei Prozent Abwertung je Jahr halbieren den Wert eines Nutzens nach etwa dreiundzwanzig Jahren --- und genau so lange dauert es, bis eine Impfung oder eine Rauchentwöhnung wirkt. Wer die Diskontierung wählt, entscheidet damit über den Rang der Vorbeugung, und das ist keine technische Frage.}
\end{minipage}
\end{center}
\end{merksatz}

\begin{merksatz}
\boxtitel{Muster 4 --- Zusammenzählen ist ein Eingriff}
\begin{center}
{\small
\renewcommand{\arraystretch}{1.48}
\begin{tabular}{@{}P{1.4cm}P{7.0cm}@{}}
\toprule
$2$ & Nutzen einer Person ist nur bis auf Streckung bestimmt --- Summen sind nicht gedeckt.\\
$10$ & Ein Verhältnis verbirgt, ob der Nenner groß oder winzig ist.\\
$13$ & Es gibt keine vernünftige Zusammenfassung mehrerer Rangfolgen.\\
$14$ & Dieselbe Summe kann gleichmäßig oder höchst ungleich verteilt sein.\\
\bottomrule
\end{tabular}}
\end{center}

\begin{center}
\begin{minipage}{0.88\textwidth}
\kursiv{\textbf{Die Rechnung für viele ist keine Fortsetzung der Rechnung für einen} --- sie setzt etwas hinzu, das Satz 2.1 nicht enthält. Man kann das tun, und man muss es tun, wenn Mittel begrenzt sind. Man sollte es nur nicht für eine bloße Verallgemeinerung halten.}
\end{minipage}
\end{center}
\end{merksatz}

\section{Was offen ist}

\begin{merksatz}
\boxtitel{Vier Fragen}
\begin{center}
{\small
\renewcommand{\arraystretch}{1.55}
\begin{tabular}{@{}P{3.0cm}P{5.6cm}@{}}
\toprule
\textbf{Wessen Werte} & Wie erhebt man Nutzen am Krankenbett, ohne die Antwort durch die Frage zu erzeugen?\\
\textbf{Rechnen und Reden} & Ob Entscheidungshilfen die Ergebnisse verbessern --- und woran man das misst.\\
\textbf{Robustheit} & Entscheiden, wenn die Wahrscheinlichkeiten selbst unsicher sind.\\
\textbf{Verteilung} & Wie man Gerechtigkeitsurteile in die Rechnung aufnimmt, ohne sie zu verstecken.\\
\bottomrule
\end{tabular}}
\end{center}

\begin{center}
\begin{minipage}{0.88\textwidth}
\kursiv{Die erste Frage ist die schwierigste. \textbf{Kapitel~12 zeigte, dass die Darstellung die Antwort mitbestimmt} --- dann ist jede Nutzenerhebung zugleich eine Beeinflussung. Es gibt kein neutrales Fragen; es gibt nur offengelegte Fragen.}
\end{minipage}
\end{center}
\end{merksatz}

\section{Sechs Sätze, die bleiben}

\begin{merksatz}
\boxtitel{Die Bilanz}
\begin{center}
{\small
\renewcommand{\arraystretch}{1.55}
\begin{tabular}{@{}P{0.7cm}P{7.7cm}@{}}
\toprule
$1$ & Eine Entscheidung besteht aus Handlungen, Zuständen und Nutzen --- und nur die ersten beiden sind Sache der Empirie.\\
$2$ & Wer vier Regeln über vernünftige Vorlieben anerkennt, maximiert einen Erwartungsnutzen.\\
$3$ & Information ist so viel wert wie die Änderung, die sie bewirkt --- und nie weniger als null.\\
$4$ & Ein Test lohnt nur zwischen zwei Schwellen; außerhalb ist die Handlung schon entschieden.\\
$5$ & Über die Zeit entscheidet man rückwärts: vom Ende her, weil jeder Schritt vom Wert des Folgenden abhängt.\\
$6$ & Aus vernünftigen Rangfolgen mehrerer Personen lässt sich keine vernünftige gemeinsame bilden --- außer diktatorisch.\\
\bottomrule
\end{tabular}}
\end{center}
\end{merksatz}

\begin{merksatz}
\boxtitel{Ein Schlusswort}
\begin{center}
\begin{minipage}{0.9\textwidth}
Dieser Band begann mit einer Patientin und einer einfachen Frage. Er endet mit derselben Frage und einer Reihe von Unterscheidungen, die vorher nicht da waren: zwischen Wahrscheinlichkeit und Bewertung, zwischen einem Befund und seiner Folge, zwischen dem Nutzen eines Menschen und der Summe über viele.

Keine dieser Unterscheidungen nimmt einem die Entscheidung ab. Sie sagen nur, worüber man sich gerade uneinig ist, und was man wissen müsste, um weiterzukommen. Das ist weniger, als die frühen Entscheidungsanalytiker hofften --- sie dachten an einen Rechenweg, der die Medizin vom Ermessen befreit. Es ist mehr, als die Skeptiker erwarteten, denn eine ausgesprochene Abwägung ist etwas anderes als eine unausgesprochene.

Am Ende steht die Einsicht, dass die schwierigsten Größen dieses Bandes nicht die Wahrscheinlichkeiten sind, sondern die Nutzen. Über jene streitet man mit Studien. Über diese muss man reden --- mit dem Patienten, im Gremium, in der Gesellschaft. Die Mathematik sagt nur, an welcher Stelle dieses Gespräch nötig ist. Dass sie das kann, ist bereits erheblich.
\end{minipage}
\end{center}

\begin{center}
\begin{minipage}{0.88\textwidth}
\kursiv{Damit ist die zweite Werkhälfte um einen weiteren Band gewachsen. Die Statistik sagt, was die Daten hergeben. Die Optimierung sagt, was das Beste ist, wenn man das Ziel kennt. Die Informationstheorie sagt, was übertragbar ist. Die Spieltheorie sagt, was übrig bleibt, wenn andere mitdenken. Das Lernen sagt, was auf den nächsten Fall zutrifft. Die Kausalität sagt, was sich ändert, wenn man handelt. \textbf{Und die Entscheidungstheorie sagt, wie man wählt, wenn man all das weiß --- und noch immer nicht genug.}}
\end{minipage}
\end{center}
\end{merksatz}
